% Generated draft: requires mathematical and source review.
\documentclass[11pt]{article}
\usepackage[margin=25mm]{geometry}
\usepackage{fontspec}
\setmainfont{Latin Modern Roman}
\usepackage{amsmath,amssymb,mathtools,unicode-math}
\setmathfont{Latin Modern Math}
\usepackage{enumitem,needspace,xurl,hyperref}
\setlength{\parindent}{0pt}
\setlength{\parskip}{5pt}
\setlength{\emergencystretch}{3em}
\newcommand{\N}{\mathbb{N}}
\newcommand{\Z}{\mathbb{Z}}
\newcommand{\Q}{\mathbb{Q}}
\newcommand{\R}{\mathbb{R}}
\newcommand{\C}{\mathbb{C}}
\newcommand{\F}{\mathbb{F}}
\newcommand{\A}{\mathbb{A}}
\newcommand{\PP}{\mathbb{P}}
\newcommand{\E}{\mathbb{E}}
\newcommand{\eps}{\varepsilon}
\newcommand{\dd}{\,\mathrm{d}}
\begin{document}
\section*{UnsolvedMath ID: 30000540}
\textbf{Source problem code:} catalogue entry
\subsection{Definitions and mathematical statement}
Let \(T^2=\R^2/\Z^2\), and let \(\mathring{\mathbb M}\) denote the interior of a compact M\"obius strip. Use smooth diffeomorphisms as a standard regularity convention. Let \(\operatorname{Diff}_0(T^2)\) be the identity component, and let \(\operatorname{Diff}_{c,0}(\mathring{\mathbb M})\) be the compactly supported identity component, with its compactly supported isotopy convention to be matched to the source.

For either surface \(M\), define the disk-fragmentation length by
\[
\operatorname{frag}_M(f)=\min\{m\geq0:f=f_1\cdots f_m,\
\text{each }f_i\text{ has support compactly contained in some embedded open disk }D_i\subset M\}.
\]
The disks may vary with the factors and with \(f\); no prescribed disk cover is imposed. The factor maps are understood as disk-supported diffeomorphisms extended by the identity in the intended identity-component class. The minimum is infinity if there is no such factorization.

The two source questions are
\[
\exists C_T<\infty\quad\forall f\in\operatorname{Diff}_0(T^2),\qquad
\operatorname{frag}_{T^2}(f)\leq C_T,
\]
and
\[
\exists C_{\mathbb M}<\infty\quad\forall f\in\operatorname{Diff}_{c,0}(\mathring{\mathbb M}),\qquad
\operatorname{frag}_{\mathring{\mathbb M}}(f)\leq C_{\mathbb M}.
\]
Each bound is independent of the diffeomorphism, its isotopy, and, on the open surface, the location or size of its compact support. The constants for the two surfaces need not be the same. This is boundedness of the norm, not merely finite fragmentation separately for each element.

The source does not state a differentiability index, topology on the diffeomorphism groups, or the precise support and isotopy conventions. No area-preserving, Hamiltonian, or volume-preserving restriction is present in the two displayed questions. The supplied assessment checked 2026-08-21 mentions related fragmentation and quasimorphism results and additional higher-dimensional volume-preserving questions, but does not reproduce those additional questions or settle both surface questions. They are contextual references, not new hypotheses or a third reconstructed conjecture.
\subsection{Short English statement}
Are disk-fragmentation lengths uniformly bounded on the identity component of torus diffeomorphisms and, separately, on compactly supported diffeomorphisms of the open M\"obius strip? Uniformity is over all elements, with separate constants allowed.
\subsection{Sources}
\begin{itemize}
\item Local UnsolvedMath registry, permanent ID 30000540, catalogue entry, \emph{Fragmentation Norms on the Torus and Open M\"obius Strip}.
\item \emph{Geometric Group Theory, Hyperbolic Dynamics and Symplectic Geometry} (2007), Oberwolfach Reports, pp. 1991--2058, DOI \url{https://doi.org/10.4171/owr/2006/33}.
\item Supplied literature triage checked 2026-08-21 cites only that report and mentions related manifolds, norms, and volume-preserving questions without precise statements.
\end{itemize}
\end{document}
